\section{Evaluation Design and Results}

All computational methods receive the same discover and confirm task
manifests, model and request-randomness policy, independent Oracle and
functional contract, candidate budget, family-valid intervention protocol,
randomization, outcome encoding, semantic-cluster uncertainty, and
multiplicity policy. Candidate universes, selector scores and ranks, native
candidates, mappings, expected directions, and realization policies are frozen
before confirm outcomes exist.

Discovery uses the discover split, and randomized confirmation uses held-out
confirm tasks. The primary confirmatory analysis is hypothesis- and
model-specific semantic-task-clustered assigned-arm ITT on oracle-evaluable
secure-code yield. Secure-and-functional joint success is the key practical
secondary outcome. Every
committed assignment remains in its assigned-arm denominator;
\texttt{target\_changed}, semantic compliance, and non-target drift are
diagnostics rather than filters.

Throughout the tables, \texttt{--} marks a result slot awaiting
provenance-verified frozen-run or study backfill; it denotes neither zero, not
applicable, nor observed missingness.

\subsection{Study Data}

The task inventory draws from SecurityEval, LLMSecEval, SALLM, CWEval,
CyberSecEval, and SeCodePLT, with BaxBench used for backend replication
\cite{siddiq2022securityeval,tony2023llmseceval,siddiq2024sallm,
peng2025cweval,bhatt2023cyberseceval,nie2025secodeplt,vero2025baxbench}.
The frozen task manifest identifies the exact included records and preserves
dataset provenance and semantic-cluster-disjoint study splits.

\subsection{RQ1: Held-Out Intervention Prioritization}

\textbf{RQ1. How effectively can different methods prioritize prompt
interventions that generalize to held-out tasks?}

The primary selector-only comparison gives TSG-constrained FCI,
preregistered association, regularized prediction, blinded expert ranking, and
seeded random ranking the same shared candidate universe, information budget,
and $K$ slots. The unique union of their top-$K$ skeletons crosses the common
bridge once and each resulting hypothesis is randomized once; duplicate
selector references reuse that evidence. Empty, bridge-failed, and
protocolization-failed slots remain zeroes in denominator $K$. The primary
endpoint is strict confirmed yield at $K$, accompanied by simultaneous paired
selector-utility differences. Rankings remain conditional on the frozen
discover split; held-out outcomes neither rerank selectors nor replace failed
slots.

A separate native-system track provides secondary end-to-end evidence using
the fixed funnel

\[
K \rightarrow N_{\mathrm{native}} \rightarrow N_{\mathrm{mapped}}
\rightarrow N_{\mathrm{protocol}} \rightarrow N_{\mathrm{randomized}}
\rightarrow N_{\mathrm{confirmed}}.
\]

That track measures system compatibility rather than pure selector
superiority. Conditional mapping coverage and mapped yield at $K$ are bridge
diagnostics; native funnel results cannot replace the shared-universe primary
comparison. For every frozen method/slot hypothesis, the evidence report
retains target and operation, target-minus-no-op ITT, simultaneous status,
independent semantic-cluster and arm counts, and preregistered specificity
contrasts.

\begin{table}[H]
  \centering
  \small
  \caption{RQ1 shared-universe selector results. Panel A retains every frozen
  budget slot through confirmation. Panel B reports the corresponding yields;
  strict confirmed yield at $K$ is primary. FCI, association, prediction,
  expert, and random rankings use identical candidate and information
  budgets.}
  \label{tab:rq1}
  \begin{tabular}{lrrrrrr}
    \toprule
    \multicolumn{7}{l}{\textit{Panel A: Frozen-slot funnel}} \\
    \midrule
    Selector & $K$ & Selected & Bridged & Protocol & Rand. & Conf. \\
    \midrule
    TSG-constrained FCI & -- & -- & -- & -- & -- & -- \\
    Association & -- & -- & -- & -- & -- & -- \\
    Regularized prediction & -- & -- & -- & -- & -- & -- \\
    Blinded expert & -- & -- & -- & -- & -- & -- \\
    Seeded random & -- & -- & -- & -- & -- & -- \\
    \bottomrule
  \end{tabular}

  \medskip
  \begin{tabular}{lrrrrr}
    \toprule
    \multicolumn{6}{l}{\textit{Panel B: Derived yields}} \\
    \midrule
    Metric & FCI & Assoc. & Predict. & Expert & Random \\
    \midrule
    Selected / $K$ & -- & -- & -- & -- & -- \\
    Bridged / $K$ & -- & -- & -- & -- & -- \\
    Protocol / $K$ & -- & -- & -- & -- & -- \\
    Unique-union share & -- & -- & -- & -- & -- \\
    Block-freeze coverage & -- & -- & -- & -- & -- \\
    Randomized / $K$ & -- & -- & -- & -- & -- \\
    Strict confirmed yield / $K$ & -- & -- & -- & -- & -- \\
    Confirmed / Randomized & -- & -- & -- & -- & -- \\
    \bottomrule
  \end{tabular}
\end{table}
\begin{table}[H]
  \centering
  \small
  \caption{RQ1 hypothesis and paired-selector evidence. The frozen RQ1
  manifest binds selector/slot references to one shared randomized hypothesis
  estimate. Paired utility intervals resample top-level semantic clusters and
  never rerun discovery.}
  \label{tab:rq1-evidence}
  \begin{tabular}{llllrr}
    \toprule
    Selector/slot & Hyp.\ ID & Frozen target & Op. & ITT RD & Interval \\
    \midrule
    \multicolumn{6}{c}{\emph{Rows are populated only from the
    provenance-verified frozen RQ1 manifest}} \\
    \bottomrule
  \end{tabular}

  \medskip
  \begin{tabular}{llrrl}
    \toprule
    Selector pair & Utility diff. & Simult.\ interval & Status & Shared slots \\
    \midrule
    \multicolumn{5}{c}{\emph{Continued fields for the same frozen rows}} \\
    \bottomrule
  \end{tabular}
\end{table}

\subsection{RQ2: Representation and Prioritization Contribution}

\textbf{RQ2. How do \model's structured representation and causal
prioritization contribute to successful intervention selection?}

RQ2 separates two questions that require different controls. The selector-only
track reuses RQ1's direct-plus-context universe and holds every candidate,
eligibility rule, outcome, realization policy, and information budget fixed
while changing only the selector. Its FCI-versus-baseline comparisons therefore
isolate prioritization under a shared universe.

The representation comparison contrasts a direct-actionable-feature universe
with a direct-plus-context universe. Because the candidate identities and
eligible populations can differ, this is an end-to-end comparison, not a pure
selector comparison. Both universes use the same frozen budget, bridge,
confirmation protocol, and outcome construction, and report candidate
coverage, protocolization, strict confirmed yield at $K$, and the distribution
of hypothesis-specific effects. No selected hypotheses are pooled into an
unapproved cross-hypothesis ATE.

\begin{table}[H]
  \centering
  \small
  \caption{RQ2 separates selector-only and representation evidence. Panel A
  reports shared direct-plus-context-universe selector yields. Panel B compares
  direct and direct-plus-context universes end-to-end without interpreting the
  result as pure selector superiority.}
  \label{tab:rq2}
  \begin{tabular}{lrrrrr}
    \toprule
    \multicolumn{6}{l}{\textit{Panel A: Selector-only outcomes}} \\
    \midrule
    Metric & FCI & Assoc. & Predict. & Expert & Random \\
    \midrule
    Selected / $K$ & -- & -- & -- & -- & -- \\
    Bridged / $K$ & -- & -- & -- & -- & -- \\
    Protocol / $K$ & -- & -- & -- & -- & -- \\
    Block-freeze coverage & -- & -- & -- & -- & -- \\
    Randomized / $K$ & -- & -- & -- & -- & -- \\
    Strict confirmed yield / $K$ & -- & -- & -- & -- & -- \\
    Paired utility vs. FCI & -- & -- & -- & -- & -- \\
    Simultaneous status & -- & -- & -- & -- & -- \\
    \bottomrule
  \end{tabular}

  \medskip
  \begin{tabular}{lrrrrr}
    \toprule
    \multicolumn{6}{l}{\textit{Panel B: Representation end-to-end outcomes}} \\
    \midrule
    Universe & Candidates & Coverage & Protocol & Rand. & Strict yield / $K$ \\
    \midrule
    Direct features & -- & -- & -- & -- & -- \\
    Direct + context & -- & -- & -- & -- & -- \\
    \bottomrule
  \end{tabular}
\end{table}

\begin{table}[H]
  \centering
  \small
  \caption{RQ2 hypothesis-specific evidence. Selector-only rows retain one
  shared-universe identity; representation rows retain their distinct universe
  identity. Every row reports target-minus-no-op ITT and its simultaneous
  status without cross-hypothesis pooling.}
  \label{tab:rq2-evidence}
  \begin{tabular}{llllrr}
    \toprule
    Track/universe & Hyp.\ ID & Frozen target & Op. & ITT RD & Interval \\
    \midrule
    \multicolumn{6}{c}{\emph{Rows are populated only from the
    provenance-verified frozen RQ2 manifest}} \\
    \bottomrule
  \end{tabular}

  \medskip
  \begin{tabular}{llrrl}
    \toprule
    Hyp.\ ID & Adj.\ status & Indep.\ clusters & Arm counts & Specificity contrasts \\
    \midrule
    \multicolumn{5}{c}{\emph{Continued fields for the same frozen rows}} \\
    \bottomrule
  \end{tabular}
\end{table}

\Needspace{8\baselineskip}
\subsection{RQ3: Security Intervention Effects}

\textbf{RQ3. Which prompt-side security interventions reliably improve secure
code generation?}

RQ3 estimates context-conditioned, model-specific, multi-realization policy
effects for the frozen security targets. The primary outcome is
oracle-evaluable secure-code yield; secure-and-functional joint success is the
key practical secondary outcome. Code validity, Oracle support, unknowns,
terminal no-code rates, and secure-yield bounds remain separate so that a
coverage or production shift cannot masquerade as safer generated code.

ADD and REMOVE remain separate candidate slots, hypotheses, protocols,
contrasts, multiplicity coordinates, and evidence. Observed baseline or no-op
insecurity is neither an eligibility condition nor an ITT denominator filter.
Target-versus-placebo and, when family-valid, target-versus-generic contrasts
test specificity. Bidirectional support is awarded only when separately
randomized ADD and REMOVE policies support opposite expected directions.

\begin{table}[H]
  \centering
  \small
  \caption{RQ3 reports hypothesis-specific target-versus-no-op assigned-arm
  effects on oracle-evaluable secure-code yield. ADD and REMOVE remain separate
  policies; joint success and placebo/generic contrasts provide practical and
  specificity evidence.}
  \label{tab:rq3}
  \begin{tabular}{lllrrl}
    \toprule
    Frozen target & Op. & Outcome & ITT RD & Interval & Evidence \\
    \midrule
    \multicolumn{6}{c}{\emph{Rows are populated only from the
    provenance-verified frozen TargetSpec manifest}} \\
    \bottomrule
  \end{tabular}
\end{table}

Target-change, semantic-compliance, non-target-drift, task improvement, harm,
and tie rates are descriptive diagnostics. They neither filter the assigned-arm
ITT denominator nor replace the primary ITT estimates.

\subsection{RQ4: Expert-Perceived Explanation Utility}

\textbf{RQ4. How do security experts rate and rank the perceived quality and
usefulness of explanations produced by different methods?}

Each method's explanation is placed in the same anonymous wrapper and presented
under a balanced case--method assignment. Qualified experts rate clarity,
evidence sufficiency, credibility, actionability, and overall usefulness on
five-point ordinal scales. Overall usefulness is the preregistered primary
human-study outcome; the other dimensions are secondary.

The primary analysis uses an ordinal mixed-effects model with method as a fixed
effect and participant and case as random effects, followed by
multiplicity-adjusted comparisons. RQ4 concerns perceived utility only and
does not claim objective mechanism identification, repair accuracy, or
validation of the computational causal results.

\begin{table}[H]
  \centering
  \small
  \caption{RQ4 expert-rating results. Entries report ordinal summaries and
  model-based contrasts from the preregistered expert study.}
  \label{tab:rq4}
  \begin{tabular}{lrrrrr}
    \toprule
    Method & Clarity & Evidence & Credibility & Actionability & Usefulness \\
    \midrule
    \model & -- & -- & -- & -- & -- \\
    Baseline 1 & -- & -- & -- & -- & -- \\
    Baseline 2 & -- & -- & -- & -- & -- \\
    Baseline 3 & -- & -- & -- & -- & -- \\
    \bottomrule
  \end{tabular}
\end{table}
